% META: DONE.1

\section{Параметрическое моделирование и история построения в САПР}

\subsection{Сущность параметрического моделирования}

Параметрическое моделирование является основным типом моделирования в современных САПР. В основе этого подхода лежит намерение выразить сложные геометрические формы изделий с помощью простых примитивов, форма которых задаётся с помощью параметров и отношений между ними (параметрами могут выступать координаты точек, величины углов, длины, радиусы окружностей или эллипсов). Отношения между параметрами называются ограничениями - для инженера-конструктора такие ограничения выглядят как размеры или геометрические отношения (параллельность, перпендикулярность, касание, концентричность и т.д.) \cite{eclass2025}.

Ключевое преимущество параметрического подхода заключается в том, что инженеру не требуется точно определять координаты всех объектов - достаточно указать геометрические взаимоотношения между ними, а всю работу по построению объекта выполнит сама программа. Это существенно сокращает затраты на внесение изменений в проект и создание новых модификаций изделий. Изменив значения параметров и заново повторив историю построения, можно получить другую геометрию, не перерисовывая модель с нуля.

\subsection{История построения как дерево операций}

История построения модели - это последовательность операций, которые были использованы для создания геометрической формы тела; и именно благодаря ей геометрическая модель становится параметрической, поскольку с каждой операцией ассоциируется набор параметров. Параметрические конструктивные элементы образуют дерево построения, которое автоматически строится в соответствии с историей создания модели и отражает взаимосвязи между элементами. В дереве 3D модели отражаются все построения и операции, присутствующие в модели. Каждая ветка дерева представляет собой историю создания одного тела модели \cite{tflexcad15}. Объекты, которые находятся глубже внутри дерева, являются более старыми, в то время как объекты, которые находятся снаружи, - более новыми, созданными на основе предыдущих.

Такая структура обладает рядом важных свойств. Во-первых, она иерархична - каждая операция может ссылаться на результаты предыдущих. Во-вторых, она упорядочена во времени - порядок операций определяет итоговую геометрию. В-третьих, она параметрически связана - изменение параметра ранней операции автоматически перестраивает все зависимые от неё элементы (если это корректно), что может позволять привносить значительные изменения в итоговую геометрию небольшими изменениями параметров одной-двух операций.

\subsection{Типовые операции и их геометрические признаки}

В типовом машиностроительном CAD дерево построения включает в себя различные типы формообразующих операций. Наиболее распространёнными являются следующие:

Операция выдавливания (Extrude/BossExtrusion) - добавление объёма путём перемещения плоского эскиза в направлении, чаще всего перпендикулярном его плоскости. Результатом является призматическое тело с поперечным сечением, соответствующим форме эскиза, а также указанной глубине выдавливания. Данная операция обладает характерными геометрическими признаками: плоские грани, параллельные плоскости эскиза, и боковые поверхности, часто ортогональные к ним.

Операция вырезания выдавливанием (CutExtrusion) - удаление объёма по тому же принципу, что и выдавливание, но в обратном направлении. Геометрический признак - вырезанная полость или отверстие, часто сквозное или заданной глубины.

Операция скругления (Fillet) - создание плавного перехода между двумя пересекающимися поверхностями путём замены острого ребра цилиндрической или тороидальной поверхностью. Геометрические признаки: плавно сопряжённые поверхности, постоянный или переменный радиус, характерное изменение кривизны.

Операция фаски (Chamfer) - снятие острого ребра под заданным углом, создание плоской наклонной грани вместо ребра. Геометрические признаки: плоская наклонная грань, углы 45° или другие заданные значения.

Каждая из этих операций оставляет на итоговой геометрии характерные «следы», которые в принципе могут быть обнаружены алгоритмически. Однако эти следы не всегда однозначны: например, отверстие может быть создано как вырезанием выдавливанием, так и операцией вращения; скругление может быть выполнено как одна операция на нескольких рёбрах или как последовательность отдельных скруглений.

\subsection{Проблема «немой» геометрии и обратный инжиниринг}

При работе с импортированными моделями (например, в формате STEP) или моделями, полученными от сторонних разработчиков, дерево построения часто отсутствует. Пользователь получает лишь итоговую геометрию - «немую» модель, из которой невозможно восстановить историю её создания. Для борьбы с этой проблемой как раз и используется обратный инжиниринг истории построения.

Как отмечается в исследованиях Autodesk, «первый набор проблем, которые мы исследовали, - это формы обратного инжиниринга. Это может включать восстановление параметрической истории из файла STEP или IGES или даже использование 3D-скана детали для построения параметрической модели». Однако, как подчёркивают те же авторы, даже на простых моделях «эскиз + выдавливание» задача далека от полного решения: «На более сложных моделях мы всё ещё не можем восстановить исходную модель за разумное время» \cite{autodesk2020}.

Таким образом, задача обратного инжиниринга конструктивной истории остаётся открытой и требует привлечения новых методов, в том числе основанных на машинном обучении.

